\section{The finite-scale contradiction}\label{sec:finish}

We now prove Theorem~\ref{thm:finite-scale}. Fix \(k\ge16\) and a
positive semidefinite Gram matrix of the corner, averaged as in
Lemma~\ref{lem:corner}. Write $\mathbf L_m$ and $\mathbf R_m$ for the matrices of
evaluated entries of $L_m$ and $\mathcal RP_+$, respectively, on this
fixed list of points. Write \(X_\ell=(\xi_\ell,\eta_\ell)\) for the nodes
\eqref{eq:nodes}; since \(b_\ell=\bar a_\ell\), \(\eta_\ell=\bar \xi_\ell\). The matrix \(\mathsf G_m\)
is PSD: each summand is an evaluated PSD kernel on the same finite list.
Its pure and mixed features may live in different spaces; their direct
sum is still a Gram. Thus \(\Phi_m\ge0\).

The crucial realigned entry is exactly a Hermitian diagonal:
\begin{equation}
 (\mathcal RP_+)(X_\ell,\overline{X_h})
 =P_+(\xi_\ell,\bar \xi_h,\bar \xi_\ell,\xi_h)
 =P_+[\xi_\ell,\bar \xi_h].                               \label{eq:realigned-diagonal}
\end{equation}
The array of such entries need not be PSD. If
\(\ell=(r,j,\sigma)\) and \(h=(r',j',\sigma')\), its coordinates for Lemma~\ref{lem:continuation} are
\[
 \alpha=(r+r')/2+i(\sigma j-\sigma'j')/2,\quad
 \beta=(\sigma j+\sigma'j')/2-i(r-r')/2.
\]
Indeed $\alpha+i\beta=a_\ell$ and $\alpha-i\beta=\bar a_h$,
which gives exactly the arguments $\xi_\ell,\bar\xi_h$ in the pure
kernel. They satisfy every bound in that lemma. Therefore
\begin{equation}
 \left|\epsilon^4 c^*\mathbf R_m c\right|
 \le\|c\|_1^2\,2^{20}\epsilon^{2^{-17}}
                    <2^{80}\epsilon^{2^{-17}}.         \label{eq:realigned-pairing}
\end{equation}

It remains to estimate the finite baseline. At the test points each
factor of $A$ and of $C$ has the form $1-e^{-\epsilon\chi}$, where
$\chi$ is a sum of two of the speeds $a_\ell,b_\ell,\bar a_h,\bar b_h$;
thus $\Re\chi\ge2$ and $|\chi|<2^{10}$. Writing
$(1-e^{-\epsilon\chi})/(\epsilon\chi)=\int_0^1e^{-\epsilon\chi u}\,du$
gives
\[
 \left|\frac{1-e^{-\epsilon \chi}}{\epsilon \chi}-1\right|
 \le\epsilon|\chi|/2\le2^{10}\epsilon.
\]
Hence \(A/\epsilon^2=A_0(1+\delta_A)\) and
\(C/\epsilon^2=C_0(1+\delta_C)\), with
\(|\delta_A|,|\delta_C|\le2^{12}\epsilon\le1/16<1/4\).
For \(|\delta|\le1/4\), use
\(|(1+\delta)^{-1}-1|\le2|\delta|\) and
\(|(1+\delta)^{-2}-1|\le6|\delta|\).
Since \(A_0,C_0\ge4\) and \(|\Delta{}|\ge1\), substitution in \eqref{eq:generating-kernels} yields
\[
 |\epsilon^4L_\infty(X_\ell,\overline{X_h})-S_F(\ell,h)|
                                               \le2^{20}\epsilon.
\]
For the bound on \(\Delta{}\), its factors are \(1+e^{-\epsilon \chi}\), and
\(|\Im(\epsilon \chi)|\le512\epsilon<1\), so the exponential has positive
real part. The estimate above follows by bounding the three reciprocal
errors by \(3\cdot2^{12}\epsilon\) in total and the remaining term by
\(4\epsilon^4\). Adding the finite tail from \eqref{eq:dyadic-tail-budget} gives, entrywise,
\begin{equation}
 |\epsilon^4L_m(X_\ell,\overline{X_h})-S_F(\ell,h)|
                                               \le2^{30}\epsilon.       \label{eq:finite-tangent-error}
\end{equation}
All slot moduli here are at most \(e^{-\epsilon}\), so Lemma~\ref{lem:tails} applies.

Finally apply \eqref{eq:correction-cancellation}, \eqref{eq:negative-pairing}, \eqref{eq:realigned-pairing}, \eqref{eq:finite-tangent-error}, and \(\|c\|_1^2<2^{60}\):
\begin{equation}
\begin{aligned}
 0\le\Phi_m
 &=\epsilon^4c^*\mathbf L_m c-\epsilon^4c^*\mathbf R_m c\\
 &<-2^{42}+2^{90}\epsilon+2^{80}\epsilon^{2^{-17}}\\
 &=-2^{42}+2^{90}m^{-1/2}+2^{80}m^{-1/262144}.
\end{aligned}                                                    \label{eq:finite-pairing-bound}
\end{equation}
This proves Theorem~\ref{thm:finite-scale}, and hence
Corollary~\ref{cor:threshold} and the negative assertion of
Theorem~\ref{thm:negative}.
